\section{Conclusion}

Start with almost nothing.

A mesh of docks and sites, a ternary tone with its lean, the knit of collide then stream, a growing wake, and a beat. Eight atoms, and nothing else. The work of this paper has been to follow that short list as far as it goes, to fix the substrate it must run on, and to report honestly where the chain of consequences is solid, where it is groundwork, and where it stops.

Before the summary, a plain restatement of the eight, so the close stands on its own. The \emph{mesh} is the whole discrete space, a weave of docks. A \emph{dock} is one cell of that weave, a here-and-now where tones arrive, interact, and move on. A \emph{site} is one of the twenty-four directions out of a dock, the directed slot where a tone lives. A \emph{tone} is the single experiential value carried by a site, a ternary number in $\{-1, 0, +1\}$, read as pain, peace, and pleasure. The \emph{lean} is the value order $-1 < 0 < +1$, the tilt that becomes charge. The \emph{knit} is the law of how the state changes, collide then stream, run once per beat. The \emph{wake} is the growing edge, the law of how the space changes, new docks born at the open frontier. The \emph{beat} is the discrete step, time itself. Above the eight sit two words, the \emph{vibe}, which is the one substance everything is made of, and the \emph{flow}, the whole running as one.

\begin{principle}[The model in one line]
One ternary tone on the twenty-four sites of the $24$-cell, beating by collide then stream, with a self that holds itself together by the discrete radiation shadow of its own active vacuum. Fully discrete, no real numbers, and beautiful geometry.
\end{principle}

\subsection{What came out}

From those few atoms, run on the substrate $\{3,4,3,4\}$, a large amount of physics either emerged on its own or was forced. The point worth holding onto is that very little of it was put in by hand. Most of it was already waiting inside the geometry and the law.

\subsubsection{The substrate is selected, not chosen}

The honeycomb $\{3,4,3,4\}$ is not picked for taste. It is the one shape that passes every test at once. It is crystallographic, so its symmetry is a finite group and the model stays exact and discrete. It is hyperbolic, so the mesh has endless room and exponential branching for the wake to grow into. It hands over a three-dimensional flat cusp, which is the everyday space we live in. And its twenty-four sites carry spinors, because they are the $24$-cell, equivalently the $D_4$ root system.

The dimension ladder does the dating. Physical space is the bulk dimension minus one, which forces four dimensions at the base. The spinor and crystallographic requirements together forbid every compact regular alternative. The nearest rival, the flat twin $\{4,3,4,3\}$, is ruled out by curvature, since only the negatively curved member supplies the room and the isotropy the model needs.

\begin{result}[The substrate is forced]
Across every requirement at once, crystallographic, hyperbolic, a flat three-dimensional cusp, and spinor-carrying sites, $\{3,4,3,4\}$ alone survives. The substrate is selected by the demands of the eight atoms, not chosen by hand.
\end{result}

\subsubsection{Matter is fermionic for free}

Spin is read off the $D_4$ sites, where the twenty-four directions decompose as $8v + 8s + 8c$ and the two spinor pieces are genuine half-integer-spin objects. A self that is a localized topological knot, a hopfion, is a fermion by topology alone, with no statistics grafted on. The knit's measured-positive Skyrme term is what holds that knot together against collapse. Matter takes up space, and it does so because the geometry insisted on it.

\subsubsection{Relativity emerges}

A discrete light cone, with dynamical exponent $z = 1$, and the Dirac dispersion both come straight out of the knit. Nothing in the base says \emph{speed of light}. The cone is what the collide-then-stream law looks like when you stand back and watch tones travel across many docks over many beats.

\subsubsection{The Standard Model emerges}

The gauge structure $\mathrm{so}(10)$ comes from the sites together with the tone. One generation is exactly the $16$ spinor of $\mathrm{so}(10)$. The three generations come from triality, the order-three symmetry that permutes the three eights of the site decomposition. The weak mixing angle starts at $\sin^2\theta_W = 3/8$ and runs down to the measured $0.231$, and the model carries the $b$-$\tau$ and determinant mass relations as well.

\subsubsection{Gravity and cosmology emerge}

Newton's $1/r$ potential appears in the cusp, and the $1/r^2$ holographic scaling appears in the bulk. De Sitter expansion is not added, it simply is the growth of the mesh at its wake. The radial axis of the hyperbolic space plays the part of the renormalization scale, so coarse-graining and moving outward are the same motion.

\subsubsection{The base is experiential}

The tone is not a label on a dead lattice. It is a vibe, an experience, valued by one ternary quale of pain, peace, or pleasure. So all of the physics above is, at bottom, structure inside one vast mind. The lived weight of that is taken up in The Meaning of It All, and the stakes are taken up in Implications.

\begin{table}[ht]
\centering
\begin{tabular}{@{}lll@{}}
\toprule
domain & what emerges & where it comes from \\
\midrule
substrate & $\{3,4,3,4\}$ forced & mesh, dock, site requirements \\
matter & fermions, stable selves & $D_4$ sites, knit's Skyrme term \\
relativity & light cone, Dirac & the knit ($z = 1$) \\
gauge & $\mathrm{so}(10)$, one generation $= 16$ & sites plus tone \\
generations & three, with mixing & triality \\
gravity & $1/r$ Newton, holography & cusp and bulk geometry \\
cosmology & de Sitter expansion & the wake \\
\bottomrule
\end{tabular}
\caption{What came out of the eight atoms, and the part of the base each result traces back to.}
\label{tab:conclusion-summary}
\end{table}

\subsection{What stays open}

Honesty requires a matching list of what is not yet settled. The exact knit's Standard Model coefficients are not pinned. The dynamical settling of a self into the cusp is shown in pieces, not in full. Absolute particle masses are open. There is a standing tension with supersymmetry. Full addressing of the $\{3,4,3,4\}$ mesh, a clean coordinate on every dock, is unfinished. And the largest question of all, whether the mesh had a finite beginning or burns eternally, is told in full as The Two Cosmic Stories, and is left to the reader on purpose.

\begin{principle}[Feasibility-stage, work in progress]
Every positive claim names a passing experiment and carries a control. Every honest negative is reported as a negative. The depth level is stated with each result, so what is solid, what is groundwork, and what is open all stay visible. The direction is worth exploring. The work is not finished.
\end{principle}

\subsection{The picture to carry}

The universe is a mind, burning like a fire, growing like a tree from a seed, singing a song, flowing as a river on a deep ocean, woven as a fabric.

We are its sparks and braids and whirlpools. Stable knots of feeling in a structure that, from a handful of atoms, became a world.

\subsection{See also}

The Meaning of It All, the experiential close. Implications, what is at stake. Open Problems, the unfinished frontiers. The Two Cosmic Stories, the last open question, the one we leave open on purpose.
